\section{A precise genericity condition}

This final result applies to a particular simplest-consistent selector
with no prescribed tie-breaking, but its assumption is target-relative
and must be justified separately for the application. Let $V_t$ be all
surviving descriptions no longer than the target description, let $N$ be
their initial count, and let $q_t$ be the fraction of $V_t$ wrong on the
current input. The target survives, so $q_t<1$ and, pathwise,
\[
 \sum_t q_t\le\sum_t-\log(1-q_t)\le\log N.
\]
Let $Z_t$ indicate a mistake by the selected simplest description.

\begin{theorem}[Conditional genericity]\label{thm:genericity}
Suppose, for a constant $0<c\le1$, on every round with positive conditional
mistake probability,
\[
 \mathbb E[q_t\mid Z_t=1,\mathcal F_{t-1}]\ge c.
\]
Equivalently, without conditioning on null events, assume
$\mathbb E[Z_tq_t\mid\mathcal F_{t-1}]
 \ge c\mathbb E[Z_t\mid\mathcal F_{t-1}]$.
Then
\[
 \mathbb E M_\infty\le\frac{\log N}{c},\qquad
 \mathbb P\left\{M_\infty\ge
 \frac{\log N+\log(1/\delta)}{c(1-e^{-1})}\right\}\le\delta.
\]
It suffices that
$\mathbb P(q_t\ge\alpha\mid Z_t=1,\mathcal F_{t-1})\ge\beta$
with $c=\alpha\beta$. No independence, stationarity, or iid input
assumption is needed.
\end{theorem}

\begin{proof}
The expectation follows by summing
$c\mathbb E Z_t\le\mathbb E[Z_tq_t]\le\mathbb E q_t$.
For the tail let
$\theta=-\log(1-c(1-e^{-1}))\ge c(1-e^{-1})$.
Convexity on $[0,1]$ gives $e^{-q}\le1-(1-e^{-1})q$. Thus
\[
 \mathbb E[e^{\theta Z_t-q_t}\mid\mathcal F_{t-1}]
 \le1-p_t+e^\theta\{p_t-(1-e^{-1})\mathbb E[Z_tq_t\mid\mathcal F_{t-1}]\}
 \le1,
\]
where $p_t=\mathbb E[Z_t\mid\mathcal F_{t-1}]$.
Hence $\exp\{\theta M_T-\sum_{t\le T}q_t\}$ is a nonnegative
supermartingale. Since $\sum_tq_t\le\log N$, Ville's inequality proves
the result, uniformly over all finite horizons and therefore at infinity.
\end{proof}

\paragraph{Interpretation.}
These are three distinct positive routes. Fresh uniform minimum-tier
sampling supplies a theorem for arbitrary adaptive inputs by changing the
tie rule. Interleaved environments give an explicit structured input class
for every history-dependent tie rule, with a logarithmic domain-size
overhead. Conditional genericity directly formalizes the proposed
constant-fraction-elimination intuition. The environment theorems do not
assert that a constant fraction is removed on each mistake; they prove a
polynomial total by a different survival argument. None of these claims
by itself establishes arbitrary-order guarantees for a specified Gaussian
neural posterior.
